\section{INTRODUCTION}

\subsection{Project Overview}
Selecting which advertisement to display to a user has become one of the most important and most-studied decision problems in the modern digital economy. Every time a user opens a content website, scrolls a social-media feed, or watches a video, a small server-side computation runs in tens of milliseconds and decides which creative, from a pool of candidates supplied by competing advertisers, will be served on that impression. This decision determines, at scale, the revenue earned by the platform, the return-on-investment realised by the advertiser, and the relevance of the content seen by the user. The decision must be made under uncertainty --- the platform never knows in advance whether a given user will click on a given advertisement --- and the only feedback it receives is whether the served creative was clicked. The information it would like to have (the click probability of every other candidate at the same moment) is permanently hidden. This is the canonical setting of the multi-armed bandit problem.

This project focuses on optimising the web-advertisement-selection decision by treating each candidate creative as an ``arm'' of a multi-armed bandit and applying reinforcement-learning algorithms to learn, online, which arm to pull. Unlike supervised-learning approaches that require a fully labelled training set up front, bandit algorithms learn continuously from the stream of impressions and clicks. They balance two competing objectives --- \emph{exploration} of new or under-served creatives in order to learn their performance, and \emph{exploitation} of known high-performing creatives to maximise immediate reward. The quality of this balance directly determines the click-through rate (CTR) and the revenue generated by the campaign.

The primary goal of the project is to develop a predictive decision system that uses five canonical bandit algorithms --- the Upper Confidence Bound (UCB1) algorithm, the $\varepsilon$-Greedy strategy with a fixed exploration rate, the Decaying $\varepsilon$-Greedy variant that anneals the exploration rate over time, the Bayesian Thompson Sampling algorithm with a Beta--Bernoulli posterior, and a Random baseline for reference. These algorithms are implemented from first principles in TypeScript and integrated into a browser-based simulation framework that runs all five agents in parallel on the same simulated ad-server, computes regret against an oracle that knows the true reward distribution, and visualises the agents' internal state and accumulated performance in real time.

Web advertising itself, sometimes called digital advertising or online advertising, is a financial arrangement that helps advertisers reach potential customers and helps content platforms generate revenue without charging users directly. It works through a marketplace in which advertisers bid on impressions, the platform serves the chosen creative, and a fee changes hands when the impression results in a click, a view, or a conversion. The size of this market is enormous: digital advertising spend exceeded six hundred billion US dollars globally in 2024, and it now accounts for the majority of all marketing spend in major economies. The extent of revenue, as well as adjacent engagement metrics such as click-through rate, dwell time, bounce rate and conversion rate, can vary widely from creative to creative and from audience to audience. A small relative improvement in selection quality, multiplied across millions of impressions a day, translates into substantial business impact.

UCB1 is a powerful, versatile, and provably efficient reinforcement-learning algorithm. It is celebrated in the bandit literature for its ability to handle the exploration--exploitation trade-off without any randomness in its selection rule and without any hyperparameter tuning beyond a single confidence constant. In the context of this project, it is deployed as the primary algorithm to optimise advertisement selection by analysing the historical click count, the running sample-mean reward, and a confidence bonus for each arm. The algorithm operates by maintaining, for every arm, an estimate of its expected reward and a confidence radius that shrinks as more samples are collected. At each step, it selects the arm whose upper confidence bound is highest, which is the arm whose true reward could most plausibly exceed the others'. Unlike static rotation, UCB1 detects even subtle differences in arm performance and concentrates impressions on the revenue-optimal creative. Unlike pure greedy methods, it never permanently abandons an arm whose early performance was unlucky --- the confidence bonus naturally re-introduces it for re-evaluation.

To improve interpretability, the simulator exposes a variety of graphical representations: cumulative-reward charts, cumulative-regret curves, rolling-CTR plots and an exploration heat-map. Each of these visualisations is tied directly to a slice of the application's live state, so the user can watch the agents make decisions step by step, watch their value estimates converge to the true CTRs, and watch the regret accumulate or flatten. When configured with realistic CTR and revenue parameters, the system is capable of producing accurate, consistent, and reliable selection decisions, making it suitable both as a research benchmark and as a teaching aid for the multi-armed bandit problem.

\subsection{Motivations}
The primary motivation behind this project is to explore and deepen our practical knowledge in the field of reinforcement learning, which is increasingly influencing advancements across many sectors --- from advertising and recommendation to robotics, healthcare, and operations research. As technology evolves rapidly, reinforcement learning has become a driving force in automating decision-making and improving efficiency in real-world applications. This project offers an opportunity to apply reinforcement-learning concepts and algorithms beyond theoretical learning by engaging in hands-on implementation of models not typically covered in classroom settings. The implementation of bandits from scratch in TypeScript, the design of a configurable environment, and the construction of a live dashboard all required us to integrate concepts from probability theory, statistics, algorithm design and front-end engineering.

One of the pressing issues in current ad-tech operations is the inefficiency of static or rule-based advertisement-selection strategies, which often waste impressions on low-performing creatives because the system does not adaptively update its serving policy. Developing a system that can adaptively learn which advertisement to show can help platforms generate more revenue, can enable advertisers to receive a fair share of impressions in proportion to their creatives' true quality, and can improve the relevance of content seen by users. Beyond the immediate technical motivation, this project gave us the opportunity to see how a well-known academic problem (the multi-armed bandit) maps cleanly onto a multi-billion-dollar industrial problem (advertisement selection at scale).

\begin{itemize}
\item \textbf{Revenue and Engagement Optimisation:} As advertising competition continues to grow, there is an urgent need for tools that help platforms, advertisers, and content owners maximise revenue and engagement effectively. A reliable bandit-based model helps platforms adjust the impression share for each creative based on observed performance, allowing for more personalised and profitable allocations. Even single-digit-percentage lifts in CTR translate, at scale, into millions of dollars of incremental revenue per quarter. Unlike A/B tests, which allocate a fixed share to each variant for the entire duration of the test, bandit policies continuously re-balance traffic toward the winner as evidence accumulates, which reduces the opportunity cost of testing.

\item \textbf{Efficient Resource Allocation:} Accurate selection decisions support advertising platforms in managing limited inventory such as page real-estate, viewer attention, and bid-budgets. By concentrating impressions on the best-performing arms, both advertiser return-on-investment and the user experience can be improved. A bandit-driven system also adapts naturally to creative fatigue: an arm that was a strong performer last month but that the audience has now tired of will see its empirical reward decay, and the policy will gradually shift impressions to fresher creatives without any human intervention.

\item \textbf{Fraud Detection and Prevention:} Reinforcement-learning systems can also play an important role in identifying suspicious patterns in click behaviour. By analysing irregularities in click rates relative to expected baselines, the system can help detect potential click fraud, reducing wasted ad spend and ultimately lowering overall costs for both advertisers and platforms. Anomalous click rates that do not translate to downstream conversions are a classic signature of bot traffic; a bandit that tracks both click and post-click signals can de-weight such arms automatically.

\item \textbf{Educational and Pedagogical Value:} The multi-armed bandit problem is foundational to modern reinforcement learning, but the standard pedagogical material is mathematical and offers little intuition about how each algorithm behaves on real reward distributions. By implementing the algorithms from scratch and exposing their internal state in a real-time dashboard, we make the exploration--exploitation trade-off observable rather than merely asserted in equations. A learner can see how a UCB1 confidence bound shrinks, how a Beta posterior sharpens, and how a poorly chosen exploration rate manifests as flat-lining regret.

\item \textbf{Reproducibility and Open Benchmarking:} A non-trivial fraction of published bandit results are difficult to reproduce because the environments, seeds and tuning details are not fully specified. Our simulator persists every run (including the random seed and the scenario configuration) to local storage, which means any reported result can be re-run bit-for-bit by the next user. This is a small but important contribution to the open-benchmarking culture in applied reinforcement learning.
\end{itemize}

\subsection{Uniqueness of the Work}
This project is unique in several respects when compared with prior bandit-simulation efforts. First, it implements all five algorithms (Random, $\varepsilon$-Greedy, Decaying $\varepsilon$-Greedy, UCB1 and Thompson Sampling) against a single unified \texttt{BaseAgent} interface, which means that comparisons are not confounded by implementation differences across libraries. Second, the simulation engine runs every selected agent \emph{synchronously} on the \emph{same} environment trajectory at every step: when the environment is asked to draw a Bernoulli outcome for arm $a$ at step $t$, the same draw is shown to every agent that picked $a$ at that step. This shared-trajectory design controls for environment noise in the comparison, which is uniquely difficult to arrange when each agent owns its own random-number generator.

Third, the environment supports two distinct reward modes --- binary click-through-rate (CTR) and continuous revenue-per-click --- with the click event drawn identically in both modes but the reward signal scaled differently. This switch is configurable from the UI, and it directly exposes one of the most operationally important lessons in applied bandit-theory: an algorithm whose modelling assumptions match the reward signal can dominate an algorithm whose assumptions do not. Thompson Sampling with a Beta--Bernoulli prior is provably near-optimal when rewards are binary clicks; under Revenue mode, its posterior is over click probability rather than revenue, and it can converge to the wrong arm. UCB1, which estimates the sample mean of the observed reward directly, sidesteps this problem.

Fourth, the system is interactive and pedagogically transparent. The internal state of every agent --- its value estimates, its current exploration rate (for $\varepsilon$-Greedy variants), its UCB scores (for UCB1), and its Beta posteriors (for Thompson Sampling) --- is rendered live in the dashboard alongside the aggregate performance metrics. This makes the algorithms' decision-by-decision behaviour inspectable, which is rare in production-grade bandit toolkits and almost unique among browser-deliverable demonstrations. Fifth, the implementation is dependency-light: no third-party reinforcement-learning library is used, which means the algorithms can be audited line-by-line from the source provided in the project repository.

Applying these techniques to a configurable synthetic environment and assessing their effectiveness with multiple evaluation metrics ensures that the comparison is both fair and dependable. The combination of from-scratch implementation, shared-trajectory benchmarking, dual-reward-mode environment design, live state visualisation, and dependency-light architecture is what distinguishes this work from off-the-shelf bandit-simulation libraries.

\subsection{Report Layout}
The report is organised to present our project in a clear and step-by-step manner, starting with an introduction that explains the significance of bandit-based advertisement selection and lays out the goals of the study. Section~2 reviews the existing literature on multi-armed bandits and ad selection, summarising the contributions of major papers from Robbins (1952) through recent contextual-bandit benchmarks, identifying the areas where this project contributes new insights, and listing the persistent gaps in current operational practice. Section~3 (\emph{Materials and Methods}) gives the details of the simulated dataset characteristics --- including reward modes, per-arm parameters, optional non-stationarity, and the encoding of categorical configuration fields. It also gives details of the choice and implementation of the five reinforcement-learning algorithms, complete with pseudocode and the update rules used in the codebase, along with the strategies used for online evaluation.

Section~4 (\emph{Results}) presents the experimental outcomes, organised first as per-agent diagnostics --- a sub-section for each algorithm with a confusion-style summary of its terminal state, its reward and regret trajectories, and its selection histogram --- and then as a comparative analysis that places all five agents on the same scale. Section~5 (\emph{Conclusions}) summarises the key findings, acknowledges the limitations of the present scope, and proposes directions for future enhancement. Section~6 lists the academic references cited in the report. Section~7 (\emph{Appendices}) gathers supplementary material: a list of abbreviations, a list of mathematical symbols, the random-number-generation and reproducibility strategy, and the threats-to-validity analysis. Section~8 contains a first-person reflection by each team member on their role and what they learned. Section~9 reproduces the originality / similarity-check summary required by the institute.
\newpage
